\documentclass[12pt]{article}
 
\usepackage{amsmath,amsthm,amssymb}

\usepackage[spanish]{babel}
\selectlanguage{spanish}
\decimalpoint
\usepackage[utf8]{inputenc}
\usepackage{amssymb,amsmath,amscd}
\usepackage{latexsym,xspace}
\usepackage{epsfig}
\usepackage{fancyhdr}
\usepackage{amsfonts}
\usepackage{mathrsfs}
\usepackage{color}
\usepackage{enumerate}
\usepackage{graphicx}
\usepackage{epstopdf}
%\usepackage[usenames,dvipsnames,svgnames,table]{xcolor}
\usepackage{xcolor}
\usepackage{amssymb,amsmath,amscd}
\usepackage{latexsym,xspace,bbm}
\usepackage[spanish]{draftcopy}
\usepackage[spanish]{babel}

\usepackage{amsfonts}
\usepackage{amsmath}
\usepackage[utf8]{inputenc}
\usepackage[spanish]{babel}
\usepackage[T1]{fontenc}
\usepackage{recycle}
\usepackage{fancyhdr}
\usepackage{graphicx}
\usepackage{enumitem}
\usepackage{float}
\usepackage{algorithm}
\usepackage{algorithmic}
\usepackage{algpseudocode}
\usepackage{amsmath}
\usepackage{tikz}
\usepackage{microtype}

\usepackage{listings}
\usepackage{xcolor}
\lstset{
    language=Python,
    backgroundcolor=\color{gray!10},
    keywordstyle=\color{blue}\bfseries,
    commentstyle=\color{green!50!black},
    stringstyle=\color{orange},
    basicstyle=\ttfamily\small,
    frame=single,
    tabsize=4,
    showstringspaces=false
}

\theoremstyle{remark}

\newcommand{\PP}{\mathbb{P}}
\newcommand{\NN}{\mathbb{N}}
\newcommand{\EE}{\mathbb{E}}
\newcommand{\RR}{\mathbb{R}}
\newcommand{\II}{\mathbb{I}}
\newcommand{\CC}{\mathbb{C}}
\newcommand{\ZZ}{\mathbb{Z}}
\newcommand{\Aa}{\mathcal{A}}
\newcommand{\Bb}{\mathcal{B}}
\newcommand{\Cc}{\mathcal{C}}
\newcommand{\Dd}{\mathcal{D}}
\newcommand{\Ee}{\mathcal{E}}
\newcommand{\Ff}{\mathcal{F}}
\newcommand{\Gg}{\mathcal{G}}
\newcommand{\Hh}{\mathcal{H}}
\newcommand{\Ii}{\mathcal{I}}
\newcommand{\Mm}{\mathcal{M}}
\newcommand{\Nn}{\mathcal{N}}
\newcommand{\Oo}{\mathcal{O}}
\newcommand{\Pp}{\mathcal{P}}
\newcommand{\Qq}{\mathcal{Q}}
\newcommand{\Uu}{\mathcal{U}}
\newcommand{\1}{\mathbbm{1}}
\newcommand{\indep}{\raisebox{0.05em}{\rotatebox[origin=c]{90}{$\models$}}}
\newtheorem{ejercicio}{Ejercicio}

% Configuración de márgenes
\addtolength{\voffset}{-1cm}
\addtolength{\hoffset}{-1.5cm}
\addtolength{\textwidth}{3cm}
\addtolength{\textheight}{2cm}

% Configuración de encabezados y pies de página
\pagestyle{fancy}
\fancyhf{}
\lhead{Profesor: Dr. Marco Vladimir Lemus Yáñez \\ Ayudantes: Naomi I. Reyes, Fernando Cruz}
\rhead{Lógica Computacional} % Cambia "ESCRIBE tu nombre" por tu nombre
\rfoot{\recycle}
\cfoot{\vspace{-0.8cm}¿Realmente necesitas imprimir esta hoja?}
\lfoot{\recycle}
\renewcommand{\headrulewidth}{1pt}
\pagenumbering{gobble}
\footskip = 50pt

% Comandos personalizados
\newcommand{\set}[1]{\left\{ #1 \right\}}

% Título del documento
\title{Práctica 2 - Punto Extra} 
\author{ESCRIBE tu nombre}
\date{\today}

\begin{document}

% Portada
\thispagestyle{empty}
\begin{figure}[ht]
  \minipage{0.76\textwidth}
  \includegraphics[width=4cm]{logo_unam.png}
  \label{EscudoUNAM}
  \endminipage
  \minipage{0.32\textwidth}
  \includegraphics[height=4.9cm, width=4cm]{logo_ciencias.png}
  \label{EscudoFC}
  \endminipage
\end{figure}

\begin{center}
  \vspace{0.1cm}
  \LARGE
  UNIVERSIDAD NACIONAL AUTÓNOMA DE MÉXICO

  \vspace{0.5cm}
  \LARGE
  FACULTAD DE CIENCIAS

  \vspace{1.5cm}
  \Large
  \textbf{Práctica 2 - Punto Extra} 

  \vspace{.01cm}
  \normalsize
  PRESENTA \\
  \vspace{.3cm}
  \large
  \textbf{Becerra Valencia César 322064287 
  \\ Cortes Nava José Luis 322115437
  \\ Díaz Anavia Javier Omar} 

  \vspace{1.3cm}
  \normalsize
  PROFESOR \\
  \vspace{.3cm}
  \large
  \textbf{Dr. Marco Vladimir Lemus Yáñez} 

  \vspace{1.3cm}
  \normalsize
  ASIGNATURA \\
  \vspace{.3cm}
  \large
  \textbf{Lógica Computacional 2026-2}

  \vspace{1.3cm}
  \today
\end{center}

    \begin{itemize}
         \item Punto extra: En pdf entregar la demostración formal que si
            \[
            \alpha \in Lpropmin \quad y \quad \beta \in Lprop
            \]
            tal que
            \[
            \alpha \equiv \beta
            \]
            entonces
            \[
            \Phi(\alpha) = \Phi(\beta)
            \] 
    \end{itemize}

    \begin{proof}
        Sea $\beta \in Lprop$ una formula proposicional y $\alpha = aMin(\beta) \in LPropmin$ su traducción. Para cualquier valor de verdad, se cumple que $\Phi(\alpha)=\Phi(\beta)$\\

        Prueba por inducción estructural sobre $\beta$\\

        \textbf{Caso Base:}    
        Sea $\beta = P_i$ un simbolo proposicional.
        Por definición, $aMin(P_i) = P_i$ y además $\Phi(P_i)=\Phi(P_i)$
        Por lo tanto $\Phi(aMin(P_i)) = \Phi(P_i)$\\

        \textbf{Hipótesis de Inducción:}
        Supongamos que para las formulas $\gamma, \delta \in Lprop$, se cumple: 
        \[
        \Phi(aMin(\gamma)) = \Phi(\gamma)
        \]
        y
        \[
        \Phi(aMin(\delta)) = \Phi(\delta) 
        \]
        
        \textbf{Paso Inductivo:}
        Probaremos entonces que la propiedad se cumple para los operadores lógicos:
        \begin{itemize}
            \item $\beta = \gamma \implies \delta$
            Por definición de $aMin$ se tiene que 
            \[
            \alpha = aMin(\gamma \implies \delta) = \neg aMin(\gamma) \lor aMin(\delta)
            \]
            Entonces
            \[
            \Phi(\neg aMin(\gamma) \lor aMin(\delta)) = H_\lor ( \Phi(\neg aMin(\gamma), aMin(\delta))) = H_\lor ( H_\neg(\Phi( aMin(\gamma)), aMin(\delta)))
            \]
            Por la Hipótesis de Inducción obtenemos:
            \[
            = H_\lor(H_\neg(\Phi(\gamma)), \Phi(\delta))
            \]
            Mientras que en $\beta$ nos queda:
            \[
            \Phi(\gamma \implies \delta) = H_{\implies} (\Phi(\gamma), \Phi(\delta)) 
            \]
            Verificaremos entonces con una tabla de verdad para ver si dan el mismo resultado:
            

                    \begin{center}
                        
                    
                    \begin{tabular}{ccccc}
                        $x$ & $y$ & $H_\neg (x)$ & $H_\lor (H_\neg (x), y)$ & $H_{\implies}(x,y) $ \\
                        F & F & V & V & V\\
                        F & V & V & V & V\\
                        V & F & F & F & F\\
                        V & V & F & V & V\\
                    \end{tabular}
\end{center}

            Como ambas columnas finales son iguales, concluimos que $$H_{\lor}(H_{\neg}(\Phi(\gamma)), \Phi(\delta)) = H_{\implies}(\Phi(\gamma), \Phi(\delta)).$$
            Por lo tanto $\Phi(\alpha) = \Phi(\beta)$.

        \end{itemize}
        

        
    \end{proof}

    
    
\end{document}
